\par\Needspace{6\baselineskip}\medskip
\begin{center}\hypertarget{native-174-article128}{}\pdfbookmark[2]{128.}{article-128}128.\end{center}
Designemus tres radios vectores eodem modo erutos ut in artt. 125, 126
per $r,r',r''$; motum angularem heliocentricum in orbita a loco secundo ad
\SourceLeaf{175}tertium per $2f$, a primo ad tertium per $2f'$, a primo ad secundum per $2f''$,
ita ut habeatur $f'=f+f''$; sit porro
\[
r'r''\sin2f=n,\qquad rr''\sin2f'=n',\qquad rr'\sin2f''=n'';
\]
denique producta quantitatis constantis $k$ (art. 2) in temporis intervalla ab
observatione secunda ad tertiam, a prima ad tertiam, a prima ad secundam
resp. $\theta,\theta',\theta''$. Incipiatur computus duplex elementorum (perinde ut in art.
praec.) tum ex $r,r',f''$ et $\theta''$, tum ex $r',r'',f,\theta$: in utroque vero calculo non
ad elementa ipsa progredieris, sed subsistes, quamprimum quantitas ea, quae
rationem sectoris elliptici ad triangulum exprimit, supraque (art. 91) per $y$ vel
$-Y$ denotata est, eruta fuerit. Sit valor huius quantitatis in calculo primo
$\eta''$, in secundo $\eta$. Habebimus itaque per formulam 18 art. 95 pro semiparametro
$p$ valorem duplicem:
\[
\BareRoot{p}=\frac{\eta''n''}{\theta''}\qquad\text{atque}\qquad
\BareRoot{p}=\frac{\eta n}{\theta}\mathpunct{\cdot}
\]
Sed per art. 82 habemus insuper valorem tertium
\[
p=\frac{4rr'r''\sin f\sin f'\sin f''}{n-n'+n''}\mathpunct{\raisebox{.45ex}{,}}
\]
qui tres valores manifesto identici esse deberent, si pro $x,y$ ab initio valores
veri accepti fuissent. Quamobrem esse deberet
\[
\frac{\theta''}{\theta}=\frac{\eta''n''}{\eta n}
\]
\[
\begin{aligned}
n-n'+n''&=\frac{4\theta\theta''rr'r''\sin f\sin f'\sin f''}{\eta\eta''nn''}\\[5pt]
&=\frac{n'\theta\theta''}{2\eta\eta''rr'r''\cos f\cos f'\cos f''}\mathpunct{\cdot}
\end{aligned}
\]
Nisi itaque his aequationibus iam in primo calculo sponte satisfit, statuere
licebit
\[
\begin{aligned}
X&=\log\frac{\eta n\theta''}{\eta''n''\theta}\\[5pt]
Y&=n-n'+n''-\frac{n'\theta\theta''}{2\eta\eta''rr'r''\cos f\cos f'\cos f''}\mathpunct{\cdot}
\end{aligned}
\]

Haec methodus applicationem aeque generalem patitur, ac secunda in
art. 125 explicata, magnum vero lucrum est, quod in hacce quinta hypotheses
primae evolutionem elementorum ipsorum non requirunt, sed in media quasi
via subsistunt. Ceterum simulatque in hac operatione eo perventum est, ut
praevideri possit, hypothesin novam a veritate haud sensibiliter discrepaturam
esse, in hac elementa ipsa vel duntaxat ex $r,r',f'',\theta''$, vel ex $r',r'',f,\theta$, vel
quod praestat ex $r,r'',f',\theta'$ determinare sufficiet.
